\documentclass[preprint,12pt]{elsarticle}

% Input the preamble file (packages and custom commands)
\usepackage{graphicx}
\usepackage{longtable}
% Layout
\usepackage[margin=1in]{geometry}

% Encoding
\usepackage[utf8]{inputenc}
\usepackage[T1]{fontenc}
\usepackage{lmodern}
\usepackage{microtype}   % Micro-typographic refinements (protrusion + expansion)
                          % to eliminate most overfull \hbox warnings automatically.
% microtype legitimately redefines \showhyphens; update the LaTeX kernel's
% stored copy so the "Command \showhyphens has changed" check passes silently.
\makeatletter
\AtBeginDocument{\let\org@showhyphens\showhyphens}
\makeatother

% Lists
\usepackage{enumitem}

% Tables
\usepackage{float}
\usepackage{rotating} % For \begin{sidewaystable}
\usepackage{booktabs} % For \toprule, \midrule, \bottomrule
\usepackage{multirow}
\usepackage{multicol}

% Symbols and Characters
\usepackage{amsmath,amssymb,amsfonts} % Added amsfonts for \mathbb
\usepackage{bm} % For bold math
\usepackage[official]{eurosym} % For € symbol
\usepackage{pifont} % For dingbats
\usepackage{newunicodechar}

% Code Listings
\usepackage{listings}

% Colors
\usepackage{xcolor}
\definecolor{codegreen}{rgb}{0,0.6,0}
\definecolor{codegray}{rgb}{0.5,0.5,0.5}
\definecolor{codepurple}{rgb}{0.58,0,0.82}
\definecolor{backcolour}{rgb}{0.95,0.95,0.92}

\lstdefinestyle{mystyle}{
    backgroundcolor=\color{backcolour},
    commentstyle=\color{codegreen},
    keywordstyle=\color{magenta},
    numberstyle=\tiny\color{codegray},
    stringstyle=\color{codepurple},
    basicstyle=\ttfamily\footnotesize,
    breakatwhitespace=false,
    breaklines=true,
    captionpos=b,
    keepspaces=true,
    numbers=left,
    numbersep=5pt,
    showspaces=false,
    showstringspaces=false,
    showtabs=false,
    tabsize=2
}
\lstset{style=mystyle}

\usepackage{ragged2e}

% Hyperlinks — must be loaded last (or near-last) to avoid option clashes.
% colorlinks colours in-text citations and cross-references; linkcolor for
% internal refs (\ref, \eqref), citecolor for \citep/\citet, urlcolor for \url.
\usepackage[colorlinks,linkcolor=blue,citecolor=blue,urlcolor=blue]{hyperref}

% Allow URLs/DOIs in the bibliography to break at any character, eliminating
% underfull \hbox warnings caused by long monospace URLs in .bbl entries.
\makeatletter
\g@addto@macro\UrlBreaks{\UrlOrds}
\makeatother
% Give TeX extra flexibility to stretch inter-word spaces in bibliography
% paragraphs, suppressing underfull \hbox warnings from long DOI/URL lines.
\emergencystretch=1em

% Suppress hyperref warnings for elsarticle internal commands (\corref, \cnotenum)
% that appear inside \begin{frontmatter} and cannot be converted to PDF strings.
\pdfstringdefDisableCommands{%
  \def\corref#1{}%
  \def\cnotenum{}%
  \def\textsubscript#1{#1}%
  \def\textsuperscript#1{#1}%
}

\begin{document}

% Input the front matter (title, authors, abstract)
\begin{frontmatter}
    \title{When Do Sustainability Priors Help Reinforcement Learning? A Hybrid FBWM-FTOPSIS-PPO Framework for Dynamic Supplier Order Allocation}

% Add your authors
% Use \author[inst1]{Name} for authors at the same institution
% Use \corref{cor1} for the corresponding author
\author[inst1]{Ali Vaezi\corref{cor1}}
\ead{ali.vaezi@studenti.polito.it}

\author[inst2]{Erfan Rabbani}
\ead{erfaninja@gmail.com}

% Add \cortext
\cortext[cor1]{Corresponding author.}

\affiliation[inst1]{
    organization={Interuniversity Department
of Regional and Urban Studies and Planning, Politecnico di Torino},
    addressline={Corso Duca degli Abruzzi, 24}, 
    city={Turin},
    postcode={10129}, 
    country={Italy}
}
\affiliation[inst2]{
    organization={Department of Management and Production Engineering, Politecnico di Torino},
    addressline={Corso Duca degli Abruzzi, 24}, 
    city={Turin},
    postcode={10129}, 
    country={Italy}
}

% Add your abstract
\begin{abstract}
Global supply chains must balance cost efficiency, sustainability, and resilience under growing demand uncertainty and disruption risk, yet existing approaches address supplier evaluation and dynamic ordering in isolation.
Static multi-criteria decision-making (MCDM) methods rank suppliers effectively but cannot adapt to changing conditions, while reinforcement learning (RL) agents learn adaptive policies but must rediscover supplier quality from scratch.
This paper proposes a two-phase framework that integrates both paradigms.
In Phase~I, the Fuzzy Best-Worst Method and Fuzzy TOPSIS translate expert judgments across 16~criteria into supplier closeness coefficients.
In Phase~II, these coefficients are embedded in the state representation of a Proximal Policy Optimization agent, creating a prior-augmented policy trained over $3\!\times\!10^{6}$~timesteps.
The framework is evaluated against a vanilla PPO baseline and a deterministic heuristic across three demand and disruption scenarios of increasing severity, totalling 27~independent runs.
Results reveal a regime-dependent trade-off: under stable conditions the prior-augmented agent achieves the lowest procurement cost (Cohen's $d = 1.59$, large effect); under systemic shocks it delivers the highest sustainability score ($d = 0.86$) with markedly lower cross-seed variance; yet under symmetric high-volatility demand, the additional state dimensions impose a learning overhead that favours the simpler baseline.
Both RL agents reduce costs by 54 to 80\,\% relative to the heuristic and maintain bullwhip ratios below unity, actively dampening demand variability.
These findings characterise MCDM-derived priors as a regime-activated asset for RL-based supply chain management.
\end{abstract}

% Add your keywords (separated by \sep)
% JCP allows 1--7 keywords. We select exactly 7, removing the generic
% "Multi-criteria decision-making" (already implied by the two specific
% method names) and "Supply chain resilience" (subsumed by "Sustainable
% supplier selection" and the title). This maximises discoverability
% while respecting the journal's guideline to avoid redundant broad terms.
\begin{keyword}
Sustainable supplier selection \sep Dynamic order allocation \sep Fuzzy Best-Worst Method \sep Fuzzy TOPSIS \sep Proximal Policy Optimization \sep Reinforcement learning \sep Bullwhip effect
\end{keyword}
\end{frontmatter}

\section{Introduction}
\label{sec:Intro}
Contemporary supply chains face compound shocks, from climatic volatility to geopolitical tension and epidemics, elevating resilience and sustainability from operational concerns to board-level priorities~\citep{Ivanov_Dolgui_2021,Li_etal_2024}. In parallel, rules such as the EU Corporate Sustainability Due Diligence Directive require end-to-end visibility across environmental, social, and governance (ESG) dimensions, complicating supplier assessment~\citep{European_Commission_2023}. Together, these forces shift supplier selection beyond price-quality trade-offs toward multi-criteria decisions that jointly consider economic performance, environmental impact, social practices, and disruption-recovery capacity under uncertainty~\citep{Dubey_etal_2021}.
By embedding sustainability and resilience evaluations directly into the ordering policy, such a framework can support cleaner procurement practices that reduce waste, lower resource consumption, and align supplier portfolios with environmental and social performance benchmarks.

Fuzzy multi-criteria decision-making (FMCDM) methods are a practical way to capture experts' linguistic judgments in sustainability evaluations without overloading respondents~\citep{Celik_Kahraman_2021,Jalali_etal_2022}. Within this family, the Fuzzy Best-Worst Method (FBWM) derives criterion weights with few pairwise comparisons and typically better consistency~\citep{Celik_Kahraman_2021}. Fuzzy TOPSIS (FTOPSIS) then ranks alternatives by their closeness to an ideal solution under fuzzy ratings~\citep{Jalali_etal_2022}. Studies in automotive sourcing, circular-economy screening, agrifood logistics, and maritime spares suggest that combining FBWM and FTOPSIS improves transparency, limits inconsistencies, and scales from SMEs to global networks~\citep{Afrasiabi_Tavana_DiCaprio_2022,El_Bettioui_Zaim_Sbihi_2023,Varchandi_Memari_AkbariJokar_2024}.

Despite these advances, FMCDM approaches are largely static: they produce point-in-time evaluations and offer limited support for dynamic order allocation as demand or disruption patterns evolve~\citep{Dubey_etal_2021}. By contrast, reinforcement learning (RL), notably Proximal Policy Optimization (PPO), shows strong adaptability in multi-echelon inventory and dynamic sourcing under stochastic conditions~\citep{Schulman_etal_2017,Wang_etal_2022}. A natural idea is therefore to embed FMCDM-derived sustainability-resilience scores directly into the RL agent's observation space, so that expert domain knowledge can guide policy learning from the first training step. Yet key questions remain unanswered: (i)~under which demand-disruption regimes do such structured priors actually improve policy quality, and under which do they impose a complexity tax that hinders generalisation? (ii)~how sensitive are the benefits to random initialisation across independent training seeds? and (iii)~do RL agents that receive MCDM priors exhibit different supply-chain dynamics, for example in bullwhip-effect dampening or supplier concentration, than agents that learn without prior information?

These questions are motivated by a Bayesian intuition: informative priors should accelerate learning when the environment activates the information they encode (e.g., supplier disruptions that correlate with the MCDM ranking), but may add noise when the environment does not discriminate among suppliers (e.g., purely demand-side volatility)~\citep{Chen_Kumar_2024}. This study tests that hypothesis empirically by training a prior-augmented PPO agent for $3\!\times\!10^{6}$~timesteps across three scenarios of increasing severity, benchmarking it against a vanilla PPO baseline and a deterministic heuristic, and replicating each configuration with three independent seeds for statistical rigour.

\subsection{Research Questions}
\begin{enumerate}
\item \textbf{RQ1}: How can FMCDM-derived supplier sustainability and resilience evaluations be systematically integrated as structured priors within a PPO-based reinforcement learning architecture to enable dynamic order allocation under uncertainty?
\item \textbf{RQ2}: Under which demand-disruption regimes (Stable Operations, High Volatility, Systemic Shock) does a prior-augmented PPO agent outperform a vanilla PPO baseline and a deterministic heuristic in terms of cost, sustainability, and bullwhip-effect dampening, and under which does the enlarged observation space impose a complexity tax?
\item \textbf{RQ3}: How do variations in random initialisation, training budget ($3\!\times\!10^{6}$~timesteps), and checkpoint selection affect the robustness and reproducibility of the integrated FBWM-FTOPSIS-PPO allocation decisions?
\end{enumerate}

\subsection{Objectives}
\begin{itemize}
\item Develop an FBWM-FTOPSIS model producing quantitative sustainability-resilience closeness coefficients suitable for RL integration.
\item Design a PPO agent whose observation space is augmented with FTOPSIS closeness coefficients and real-time disruption flags, and compare it against a vanilla PPO baseline and a deterministic base-stock heuristic across three demand-disruption regimes.
\item Validate adaptivity, robustness, and regime-dependent trade-offs via multi-seed replication (27~experimental runs), $3\!\times\!10^{6}$-step training with checkpoint-based safety mechanisms, and comprehensive statistical analysis (Welch's $t$-tests, Cohen's~$d$, one-way ANOVA).
\end{itemize}

\section{Literature Review}

\subsection{Sustainable and Resilient Supply Chains}
Sustainable and resilient supply chains have become strategic imperatives as disruptions from climate change, geopolitical instability, and health crises intensify. Current approaches aim not only to reduce environmental footprints and social risks but also to secure rapid recovery and continuity when shocks occur~\citep{Ivanov_Dolgui_2021}. Regulatory pressure, most notably the EU Corporate Sustainability Due Diligence Directive, now requires firms to monitor upstream impacts across ESG dimensions, turning supplier selection into a strategic, multi-stakeholder decision~\citep{European_Commission_2023}. Resilience research has moved from conceptual taxonomies such as agility, flexibility, robustness, and visibility to quantitative metrics embedded in decision models, underscoring the need for real-time adaptation and continuous learning under uncertainty~\citep{Dubey_etal_2021}.

\subsection{Fuzzy MCDM Methods for Supplier Evaluation}

\subsubsection{Fuzzy Best-Worst Method (FBWM)}
The Fuzzy Best-Worst Method (FBWM) extends classic BWM by representing expert judgments with fuzzy triangular numbers, which captures linguistic vagueness in pairwise comparisons and keeps elicitation effort low through few comparisons~\citep{Celik_Kahraman_2021}. Empirical studies report higher consistency ratios and lower inconsistency indices than AHP and crisp BWM, particularly in sustainability trade-off settings~\citep{haseli_2024_sf_bwm_cis,dong_etal_2021_gscm}. Nonetheless, FBWM is static: it provides point-in-time weights and does not adapt when conditions change.

\subsubsection{Fuzzy TOPSIS (FTOPSIS)}
Fuzzy TOPSIS (FTOPSIS) ranks alternatives by their relative closeness to an ideal solution in a fuzzy decision space, allowing simultaneous consideration of best and worst case scenarios under imprecise data~\citep{Jalali_etal_2022}. When combined with FBWM-derived weights, FTOPSIS has produced transparent supplier rankings in manufacturing, food processing, and services, and has been reported to outperform entropy-based and probabilistic approaches in robustness tests~\citep{Zhang_Zong_Dou_Zhao_Tang_Li_2021}. However, like FBWM, FTOPSIS is static; it yields point-in-time rankings and does not adapt as demand and disruption regimes evolve.

\subsection{Synergies, Recent Trends, and Advanced Enhancements}
Hybrid FBWM-FTOPSIS frameworks have evolved along two enhancement directions.

% Trimmed: The original two paragraphs detailed specific fuzzy extensions
% (Pythagorean, q-rung orthopair) and individual ML-MCDM studies at length.
% The revision condenses both into a single paragraph per direction, removing
% implementation details that are tangential to this study's contribution.

First, advanced fuzzy modelling, including enhanced defuzzification techniques and extended fuzzy environments such as Pythagorean and q-rung orthopair sets, has improved ranking stability and sensitivity under uncertainty, particularly in multi-stakeholder sustainability contexts~\citep{Recent_fuzzy_extensions_2024}.

Second, machine learning (ML) mechanisms are increasingly integrated into MCDM architectures.
Hybrid ML-MCDM frameworks employ dimensionality reduction, predictive modelling, and data-driven weighting to reduce subjectivity in weight elicitation and improve analytical scalability~\citep{Recent_ML_MCDM_hybrid_2024}.
Two-phase models combining demand-forecasting techniques with fuzzy weighting and multi-objective programming demonstrate how predictive accuracy influences supplier selection outcomes~\citep{Recent_forecast_MCDM_2025}, while dimensionality-reduction-assisted approaches address high-dimensional evaluation structures~\citep{Recent_dimreduction_MCDM_2025}.

Despite improved robustness and scalability, most enhanced frameworks remain embedded within static evaluation paradigms.
ML components are typically applied for preprocessing or weight derivation, while final rankings remain point-in-time outputs.
Integration into adaptive decision-making architectures capable of continuous learning under dynamic disruption regimes remains limited.

\subsection{Review of Recent Related Works}
Recent research extends beyond static hybridisation by exploring the integration of FMCDM with reinforcement learning (RL) to enable adaptive supplier and inventory decisions. \citet{Wang_etal_2022} embed crisp MCDM weights into a Q-learning framework for multi-echelon inventory control and report improved service levels under stochastic demand. Other studies incorporate AHP-TOPSIS scores as initial preference structures within Deep Q-Network (DQN) agents for order allocation, observing faster convergence but increased sensitivity to exploration hyperparameters~\citep{Recent_MCDM_DQN_2024}.

Proximal Policy Optimization (PPO) has gained prominence in supply-chain control due to its stability, sample efficiency, and robustness in continuous decision spaces~\citep{Schulman_etal_2017}. Recent applications include PPO with adaptive entropy tuning for real-time replenishment under demand uncertainty~\citep{Li_etal_2024}, demonstrating improved adaptability relative to fixed-entropy configurations. However, these PPO-based studies typically operate without structured domain knowledge: the agent must discover supplier quality purely from reward feedback, which can slow convergence and increase sensitivity to random initialisation.

Despite these advances, existing RL integrations typically rely on crisp MCDM outputs or heuristic initialisation strategies. Structured incorporation of fuzzy-derived sustainability-resilience evaluations within the observation space of policy-gradient architectures, particularly PPO-based agents, remains underexplored.

% Removed: A paragraph summarising sector-specific FBWM-FTOPSIS applications
% (automotive, circular-economy, agrifood, maritime) was deleted because these
% studies were already cited in the Introduction (Section 1, paragraph 2) and
% added no new analytical insight beyond confirming cross-sector robustness.

\subsection{Identified Gaps}
Despite these advances, critical gaps impede seamless integration of sustainability-resilience evaluation with dynamic order allocation:
\begin{enumerate}
    \item \textbf{Static evaluations.} FBWM-FTOPSIS frameworks produce point-in-time weights and rankings, limiting responsiveness to demand volatility and evolving disruption regimes~\citep{Dubey_etal_2021}.
    \item \textbf{Lack of structured prior integration.} Expert-derived sustainability-resilience scores have not been effectively incorporated into the observation space of RL agents, missing an opportunity to preserve domain knowledge during policy learning~\citep{Chen_Kumar_2024}.
    \item \textbf{Scenario-dependent evaluation gap.} Existing RL studies in supply chain management seldom evaluate the same agent architecture across multiple disruption regimes (stable, volatile, and systemic-shock conditions), making it difficult to characterise when domain priors help and when they may hinder learning~\citep{Schulman_etal_2017}.
    \item \textbf{Seed sensitivity and reproducibility.} Deep RL results in operations research are frequently reported from a single random seed, obscuring the inter-run variability that is critical for managerial decision-making under uncertainty~\citep{Li_etal_2024}.
\end{enumerate}

\subsection{Positioning and Contribution of the Current Study}
This study addresses these gaps with a two-phase hybrid framework: (1)~FBWM-FTOPSIS to derive comprehensive sustainability-resilience closeness coefficients for a pool of five suppliers, and (2)~embedding these coefficients directly into the observation space of a PPO agent so that the MCDM scores function as structured priors during policy learning. The proposed prior-augmented agent (M\textsubscript{3}) is benchmarked against a vanilla PPO baseline without prior information (M\textsubscript{2}) and a deterministic base-stock heuristic (M\textsubscript{1}) across three demand-disruption scenarios of increasing severity: Stable Operations, High Volatility, and Systemic Shock. Each configuration is replicated with three independent random seeds to quantify inter-run variability and provide statistical rigour. Through this multi-scenario, multi-seed design, we evaluate the framework's adaptivity, robustness, and transparency in dynamic order allocation.

\section{Methodology}
\label{sec:methodology}
\subsection{Problem Description}
% Trimmed: The original opening restated the limitations of static FMCDM and the
% promise of RL already covered in Sections 1--2, adding ~100 words of redundancy.
% The revision jumps directly to the proposed framework.
We consider a single buyer sourcing two products from five heterogeneous suppliers under non-stationary demand and stochastic disruptions.
The buyer's objective is to minimise total operational cost while maintaining a high sustainability-resilience allocation profile.
To bridge static expert evaluation and adaptive ordering, we propose a two-phase hybrid framework.
First, sustainability and resilience closeness coefficients generated by FBWM-FTOPSIS serve as structured priors that are appended to a PPO agent's observation vector, guiding it toward high-value supplier options from the outset.
The resulting prior-augmented agent (M\textsubscript{3}, 16-dimensional observation space) is compared with a vanilla PPO baseline that receives only inventory, demand, and pipeline information (M\textsubscript{2}, 6-dimensional observation space) and a deterministic base-stock heuristic (M\textsubscript{1}).
All three models are trained and evaluated across three demand-disruption scenarios (Stable Operations, High Volatility, and Systemic Shock) with three independent random seeds per configuration, yielding 27 experimental runs.

\subsection{Criterion System Design}
% Trimmed: Each dimension paragraph was condensed to a brief overview
% sentence plus one-line criterion descriptions, since Table 1 already
% lists all 16 criteria. Detailed justifications were removed to save
% approximately 400 words.
To assess suppliers across sustainability and resilience, this study uses a four-category, 16-criterion set (Table~\ref{tab:criteria}) selected from prior literature, expert interviews, and sectoral relevance.

\paragraph{Economic criteria.}
Cost competitiveness, product quality, delivery performance, and financial capability capture traditional supplier performance and operational feasibility~\citep{Afrasiabi_Tavana_DiCaprio_2022,christopher_peck_2004_resilient_sc}.

\paragraph{Environmental criteria.}
Critical-minerals stewardship, eco-labels and certifications (e.g., ISO~14001), reuse and recycling initiatives, and waste management and pollution control assess ecological impact and circular-economy alignment~\citep{Varchandi_Memari_AkbariJokar_2024,testa_etal_2012_gscm_drivers}.

\paragraph{Social criteria.}
Occupational health and safety, social compliance and grievance mechanisms, labour-contract formalisation, and employment insurance and benefits ensure ethical practices and stakeholder well-being~\citep{gualandris_kalchschmidt_2016_sustainability_trust,El_Bettioui_Zaim_Sbihi_2023}.

\paragraph{Resilience criteria.}
Agility, flexibility, robustness, and visibility and risk awareness evaluate a supplier's capacity to detect, absorb, and recover from disruptions~\citep{soni_etal_2014_sc_resilience,pettit_etal_2010_scr_framework,brandonjones_etal_2014_scr_robustness,Ivanov_Dolgui_2021}.

\begin{longtable}{@{}ll@{}}
\caption{Overview of Evaluation Criteria by Sustainability Dimension}\label{tab:criteria}\\
\toprule
Dimension & Criterion \\
\midrule
\endfirsthead
\toprule
Dimension & Criterion \\
\midrule
\endhead
Economic & Cost Competitiveness \\
Economic & Product Quality \\
Economic & Delivery Performance \\
Economic & Financial Capability \\
Environmental & Critical-minerals sourcing \& recycled content \\
Environmental & Eco-Labeling and Certification \\
Environmental & Reuse and Recycling \\
Environmental & Waste Management and Pollution Control \\
Social & Occupational Health and Safety \\
Social & Social compliance \& grievance systems \\
Social & Labor Contract Formalisation \\
Social & Employment Insurance and Benefits \\
Resilience & Agility \\
Resilience & Flexibility \\
Resilience & Robustness \\
Resilience & Visibility and Risk Awareness \\
\bottomrule
\end{longtable}

\subsection{Phase I, part I: Fuzzy Best-Worst Method (FBWM) for Criterion Weight Determination}

\subsubsection{Theoretical foundation and integration}
% Trimmed: The original two paragraphs restated FBWM's advantages (fewer
% comparisons, linguistic encoding) already covered in Section 2.2.1.
% The revision retains only the procedural rationale for using FBWM in Phase I.
FBWM is selected as the weighting component because its $2n{-}3$ pairwise comparisons (versus $n(n{-}1)/2$ in AHP) reduce elicitation burden while maintaining consistency~\citep{haseli_2024_sf_bwm_cis,Celik_Kahraman_2021}. Linguistic assessments are encoded as triangular fuzzy numbers, and the resulting fuzzy weights feed directly into the FTOPSIS ranking phase, preserving computational consistency across the integrated framework~\citep{petrovic_etal_2022_axioms_efqm_fbwm_fmoora}.

\subsubsection{Mathematical framework and fuzzy-number operations}
FBWM represents imprecise judgments with triangular fuzzy numbers (TFNs), where each TFN is $\tilde{A}=(l,m,u)$ with $l$, $m$, and $u$ denoting the lower bound, modal value, and upper bound, respectively~\citep{li_etal_2021_sc_fintech,haseli_2024_sf_bwm_cis}. For positive TFNs, standard fuzzy arithmetic applies~\citep{El_Bettioui_Zaim_Sbihi_2023,li_etal_2021_sc_fintech}:
\[
\text{Addition: } \tilde{A}_1 \oplus \tilde{A}_2 = (l_1+l_2,\; m_1+m_2,\; u_1+u_2)
\]
\[
\text{Multiplication: } \tilde{A}_1 \otimes \tilde{A}_2 = (l_1 \times l_2,\; m_1 \times m_2,\; u_1 \times u_2)
\]
\[
\text{Division: } \tilde{A}_1 \oslash \tilde{A}_2 = (l_1/u_2,\; m_1/m_2,\; u_1/l_2)
\]
A seven-point linguistic scale maps verbal judgments to TFNs, for example, from equally important $(1,1,1)$ to very important $(2.5,3,3.5)$, supporting intuitive expert input while preserving mathematical tractability~\citep{ratandhara_kumar_2023_bwm,El_Bettioui_Zaim_Sbihi_2023}.

\subsubsection{Best-to-Others and Others-to-Worst comparisons}
The FBWM procedure starts by asking experts to identify the best and worst criteria among the 16 sustainability-resilience criteria in Table~2~\citep{Afrasiabi_Tavana_DiCaprio_2022,gan_2019_resilient_supplier_ddns}. Using linguistic terms, they then provide fuzzy pairwise preferences that form two vectors: \\
\textbf{Best-to-Others (BO) vector:} $\tilde{A}_B=(\tilde{a}_{B1},\tilde{a}_{B2},\ldots,\tilde{a}_{Bn})$, where $\tilde{a}_{Bj}$ is the fuzzy preference intensity of the best criterion over criterion $j$~\citep{El_Bettioui_Zaim_Sbihi_2023}. \\
\textbf{Others-to-Worst (OW) vector:} $\tilde{A}_W=(\tilde{a}_{1W},\tilde{a}_{2W},\ldots,\tilde{a}_{nW})$, where $\tilde{a}_{jW}$ is the fuzzy preference of criterion $j$ over the worst criterion~\citep{haseli_2024_sf_bwm_cis,ratandhara_kumar_2023_bwm}. \\
Together, these vectors capture uncertainty in expert judgments while providing enough structure to compute consistent criterion weights via optimisation~\citep{li_etal_2021_sc_fintech,petrovic_etal_2022_axioms_efqm_fbwm_fmoora}.

\subsubsection{Group aggregation and weight optimisation}
In group settings with $k$ experts, each expert supplies BO/OW vectors; these are aggregated with the geometric mean, which preserves the multiplicative-reciprocal property needed for BWM consistency~\citep{wanmohd_abdullah_2016_supplier_selection}:
\[
\tilde{a}_{Bj} \;=\; \Bigg(\prod_{p=1}^{k} \big(\tilde{a}_{Bj}^{(p)}\big)\Bigg)^{1/k}
\]
The optimal fuzzy weights $\{\tilde{w}_j\}$ are then obtained by solving a nonlinear program that minimises the maximum deviation between stated preferences and implied weight ratios~\citep{rezaei_2019_nonlinear_bwm,petrovic_etal_2022_axioms_efqm_fbwm_fmoora}:
\[
\text{Minimise: } \tilde{\xi}=(\xi,\xi,\xi)
\]
\[
\text{Subject to: } 
\begin{cases}
\big|\, \tilde{w}_B \oslash \tilde{w}_j \ominus \tilde{a}_{Bj} \,\big| \le \tilde{\xi}, & \forall j \\
\big|\, \tilde{w}_j \oslash \tilde{w}_W \ominus \tilde{a}_{jW} \,\big| \le \tilde{\xi}, & \forall j \\
\sum_{j} \mathrm{GMIR}(\tilde{w}_j) = 1, \\
\tilde{w}_j \succeq 0, \quad \forall j
\end{cases}
\]
where GMIR denotes the graded mean integration representation used for defuzzification~\citep{li_etal_2021_sc_fintech,petrovic_etal_2022_axioms_efqm_fbwm_fmoora}.

\subsubsection{Solution procedure (crisp LP decomposition)}
In practice, the fuzzy nonlinear model is solved robustly by decomposing it into a series of crisp linear programs, which improves tractability and numerical stability. We follow the approach of \citet{guo_zhao_2017_fbwm}: solve three linear programs, one for each TFN component (lower, modal, upper). Each program minimises the consistency deviation $\xi$ under the linearized constraints
\[
|\, w_B - a_{Bj}\, w_j \,| \le \xi\, w_j 
\quad \text{and} \quad
|\, w_j - a_{jW}\, w_W \,| \le \xi\, w_W.
\]
The three resulting weight vectors $(w^l,w^m,w^u)$ are then combined to form the fuzzy weights $\tilde{w}_j=(l_j,m_j,u_j)$, and sorted so that the triangular property $l_j \le m_j \le u_j$ holds.

\subsubsection{Consistency verification and transition to ranking}
Consistency of expert judgments is checked via the fuzzy consistency ratio,
\[
CR=\xi^*/CI,
\]
where $\xi^*$ is the optimal objective value and $CI$ is the consistency index tied to the fuzzy best-to-worst preference value~\citep{rezaei_2019_nonlinear_bwm,guo_zhao_2017_fbwm}. We adopt the common threshold $CR \le 0.1$ as acceptable, ensuring reliable weight derivation~\citep{haseli_2024_sf_bwm_cis}.
Once consistency is satisfactory, the fuzzy criterion weights $\tilde{w}_j$ are normalized and passed directly to the FTOPSIS weighted evaluation matrix, forming the bridge from weighting to ranking~\citep{Afrasiabi_Tavana_DiCaprio_2022,petrovic_etal_2022_axioms_efqm_fbwm_fmoora}. This seamless handoff preserves expert-derived sustainability and resilience priorities and operationalizes them in supplier evaluation, maintaining coherence across the full framework.

\subsection{Phase I, part II: Fuzzy TOPSIS (FTOPSIS) for Sustainable Supplier Ranking}

\subsubsection{Theoretical Foundation and Methodological Integration}
% Trimmed: The original two paragraphs restated FTOPSIS's dual-objective
% motivation and its compatibility with FBWM, already discussed in
% Section 2.2.2. The revision retains only the essential procedural link.
FTOPSIS ranks alternatives by their fuzzy closeness to an ideal solution, using the same TFN structure as FBWM to maintain mathematical consistency~\citep{saleem_etal_2025_fuzzy_topsis_gscm,zare_etal_2023_pf_bwm_supplier}. Its dual objective (maximising proximity to the ideal, minimising distance from the anti-ideal) supports balanced trade-offs across conflicting sustainability and resilience criteria~\citep{amiri_aref_etal_2012_new_fuzzy_pnis_topsis}.

\subsubsection{Fuzzy Decision Matrix Construction and Normalisation}
Consistent with FBWM (Section 3.3), FTOPSIS uses the same triangular fuzzy number structure $\tilde{A}=(l,m,u)$ and the same linguistic scale to keep the framework coherent~\citep{zare_etal_2023_pf_bwm_supplier}.
For group decisions with $k$ experts, individual supplier ratings are aggregated by the geometric mean:
\[
\tilde{x}_{ij}=\left(\prod_{p=1}^{k} x_{ij}^{(p)}\right)^{1/k},
\]
where $x_{ij}^{(p)}$ is expert $p$'s fuzzy rating for supplier $i$ on criterion $j$~\citep{saghafian_hejazi_2005_modified_fuzzy_topsis,zare_etal_2023_pf_bwm_supplier}.

\subsubsection{Fuzzy decision matrix and normalisation}
The aggregated ratings form the fuzzy decision matrix $\tilde{D}=[\tilde{x}_{ij}]_{m\times n}$, which is then linearly normalised to make heterogeneous criteria comparable~\citep{sarraf_mcguire_2021_normalization_topsis}:
\[
\text{Benefit criteria:}\quad
\tilde{r}_{ij}=\left(\frac{a_{ij}}{c_j^*},\, \frac{b_{ij}}{c_j^*},\, \frac{c_{ij}}{c_j^*}\right)
\]
\[
\text{Cost criteria:}\quad
\tilde{r}_{ij}=\left(\frac{a_j^-}{c_{ij}},\, \frac{a_j^-}{b_{ij}},\, \frac{a_j^-}{a_{ij}}\right)
\]
Here, $c_j^*=\max\{c_{ij}\}$ and $a_j^-=\min\{a_{ij}\}$ denote, respectively, the maximum and minimum values observed for criterion $j$ across all suppliers~\citep{sarraf_mcguire_2021_normalization_topsis}. 

\subsubsection{Weighted Normalisation and Ideal-Solution Determination}
The normalized fuzzy decision matrix is weighted using the FBWM-derived fuzzy criterion weights $\tilde{w}_j$, creating a direct link between the two phases~\citep{zare_etal_2023_pf_bwm_supplier}:
\[
\tilde{v}_{ij}=\tilde{r}_{ij}\otimes \tilde{w}_j.
\]
Next, the Fuzzy Positive Ideal Solution (FPIS) and Fuzzy Negative Ideal Solution (FNIS) are defined as~\citep{saghafian_hejazi_2005_modified_fuzzy_topsis,amiri_aref_etal_2012_new_fuzzy_pnis_topsis}:
\[
A^{+}=\left\{\max\{\tilde{v}_{ij}\mid i=1,\ldots,m\}\,\middle|\, j\in B;\;
\min\{\tilde{v}_{ij}\mid i=1,\ldots,m\}\,\middle|\, j\in C\right\},
\]
\[
A^{-}=\left\{\min\{\tilde{v}_{ij}\mid i=1,\ldots,m\}\,\middle|\, j\in B;\;
\max\{\tilde{v}_{ij}\mid i=1,\ldots,m\}\,\middle|\, j\in C\right\},
\]
where $B$ and $C$ denote the sets of benefit and cost criteria, respectively~\citep{collan_etal_2015_fpom_rnd}.

\subsubsection{Distance Calculation and Ranking Determination}
Using the vertex method, compute the distances between each supplier's weighted performance and the ideal solutions~\citep{chen_2000_fuzzy_topsis,sarraf_mcguire_2021_normalization_topsis}:
\[
d_i^{+}=\sum_{j=1}^{n} d\!\left(\tilde{v}_{ij},\, \tilde{v}_{j}^{+}\right),\qquad
d_i^{-}=\sum_{j=1}^{n} d\!\left(\tilde{v}_{ij},\, \tilde{v}_{j}^{-}\right),
\]
where, for two triangular fuzzy numbers $\tilde{A}=(a_1,a_2,a_3)$ and $\tilde{B}=(b_1,b_2,b_3)$,
\[
d(\tilde{A},\tilde{B})=\sqrt{\frac{(a_1-b_1)^2+(a_2-b_2)^2+(a_3-b_3)^2}{3}} \quad \text{\citep{chen_2000_fuzzy_topsis,izadikhah_2009_hamming_fuzzy_topsis}}.
\]
The closeness coefficient then determines the final ranking~\citep{collan_etal_2015_fpom_rnd}:
\[
CC_i=\frac{d_i^{-}}{d_i^{+}+d_i^{-}}.
\]
Suppliers are ranked in descending order of $CC_i$; higher coefficients indicate stronger sustainability-resilience performance~\citep{collan_etal_2015_fpom_rnd,zare_etal_2023_pf_bwm_supplier}.

\subsubsection{Validation and Sensitivity Analysis}
The FBWM-FTOPSIS integration is validated via the Phase~I consistency ratio check and a sensitivity analysis on criterion-weight perturbations~\citep{raveena_umamaheswari_2025_pharma_supplier_mcdm,chen_etal_2021_fuzzy_topsis_supplier}.
Robustness is assessed by testing ranking stability under $\pm 20\%$ weight variations, and a correlation analysis between fuzzy weights and supplier scores verifies methodological coherence~\citep{rostami_etal_2023_gp_fbwm_viable_ss,liou_etal_2021_mcdm_datamining_green_supplier}.
Together, these checks provide reliable decision support for dynamic, sustainable supply-chain settings where both expert-judgment quality and ranking stability are critical~\citep{chaharsooghi_ashrafi_2014_neofuzzy_topsis,chen_etal_2021_fuzzy_topsis_supplier}.

\subsection{Phase II: Reinforcement Learning for Dynamic Order Allocation}
Phase~I yields a static ranking of the five suppliers via FTOPSIS closeness coefficients $CC_i \in [0,1]$. Phase~II operationalises these scores within a reinforcement-learning loop, enabling the agent to adjust order quantities daily in response to evolving demand, lead-time realisations, and supplier disruptions.

\subsubsection{Markov Decision Process Formulation}
The supply-chain environment is modelled as an episodic Markov Decision Process $(\mathcal{S},\mathcal{A},P,r,\gamma)$ implemented in the OpenAI Gymnasium framework. Each episode spans $T=365$ time steps, corresponding to one year of daily operations with $n=5$ heterogeneous suppliers and $m=2$ products.

\paragraph{State space.}
The full state vector $\mathbf{s}_t\in\mathbb{R}^{16}$ comprises four groups:
\begin{equation}
\mathbf{s}_t =
\bigl[\,
\underbrace{I_{t,1},\,I_{t,2}}_{\text{inventory}},\;
\underbrace{d_{t-1,1},\,d_{t-1,2}}_{\text{last demand}},\;
\underbrace{\Pi_{t,1},\,\Pi_{t,2}}_{\text{pipeline}},\;
\underbrace{\delta_{t,1},\ldots,\delta_{t,5}}_{\text{disruption flags}},\;
\underbrace{CC_1,\ldots,CC_5}_{\text{FTOPSIS scores}}
\,\bigr],
\label{eq:state}
\end{equation}
where $I_{t,u}$ is the on-hand inventory of product~$u$, $d_{t-1,u}$ is the previous day's realised demand, $\Pi_{t,u}$ is the total pipeline (in-transit) inventory, $\delta_{t,i}\in\{0,1\}$ is the disruption flag for supplier~$i$, and $CC_i$ is the time-invariant FTOPSIS closeness coefficient from Phase~I.

\paragraph{Action space.}
The agent outputs a continuous action $\mathbf{a}_t\in[-1,1]^{5\times 2}$ that is rescaled to physical order quantities:
\begin{equation}
q_{t,i,u} = \frac{a_{t,i,u}+1}{2}\;\cdot\;\bar{Q}_{i,u},
\label{eq:action}
\end{equation}
where $\bar{Q}_{i,u}$ is the maximum daily capacity of supplier~$i$ for product~$u$. This normalised representation improves numerical stability and encourages balanced exploration by the policy gradient optimiser.

\paragraph{Demand model.}
Daily demand follows a non-stationary composite process:
\begin{equation}
d_{t,u} = \max\!\bigl(0,\;\mu_u + \alpha\,t + A\sin(2\pi t/T) + \xi_t + \varepsilon_{t,u}\bigr),
\label{eq:demand}
\end{equation}
where $\mu_u$ is the base demand, $\alpha$ the linear trend coefficient, $A$ the seasonality amplitude, $\xi_t$ is a demand-shock term that fires with probability $p_{\text{shock}}$, and $\varepsilon_{t,u}\sim\mathcal{N}(0,\sigma_u^2)$ is Gaussian noise. Scenario-specific parameterisations of $\sigma_u$ and $p_{\text{shock}}$ are given in Table~\ref{tab:scenarios}.

\paragraph{Reward function.}
The reward balances operational cost minimisation with sustainability:
\begin{equation}
r_t = -\hat{C}_{t}^{\text{ops}} + \lambda_{\text{sust}}\,\hat{F}_t,
\label{eq:reward}
\end{equation}
where $\hat{C}_{t}^{\text{ops}} = (c_t^{\text{purchase}}+c_t^{\text{holding}}+c_t^{\text{shortage}})/C_{\text{ref}}$ is the daily operational cost normalised by a reference daily cost $C_{\text{ref}}=\text{\euro}\,10{,}000$, and $\hat{F}_t = \sum_i q_{t,i}\,CC_i\;/\;\sum_i q_{t,i}$ is the order-weighted mean FTOPSIS score. The sustainability weight is set to $\lambda_{\text{sust}}=0.1$ throughout the main experiments.

\subsubsection{Agent Architectures}
Three models are compared; all share the same environment and cost structure to isolate the effect of prior information on policy quality.

\paragraph{M\textsubscript{1}: Base-Stock Heuristic.}
A deterministic inventory-management heuristic that orders exclusively from the highest-scoring supplier ($\arg\max_i CC_i$). The order quantity satisfies
$Q_t = \max(0,\,S - I_t - \Pi_t)$,
where $S$ is a fixed base-stock target calibrated to approximately four days of average demand. M\textsubscript{1} serves as a non-learning domain-expert baseline.

\paragraph{M\textsubscript{2}: Vanilla PPO.}
A standard Proximal Policy Optimisation agent~\citep{Schulman_etal_2017} that receives a \emph{reduced} 6-dimensional observation $\mathbf{s}_t^{(6)} = [I_{t,1},I_{t,2},d_{t-1,1},d_{t-1,2},\Pi_{t,1},\Pi_{t,2}]$ via a wrapper that strips the disruption flags and FTOPSIS scores from the full state. M\textsubscript{2} must discover supplier quality purely through reward feedback.

\paragraph{M\textsubscript{3}: PPO + FTOPSIS Priors.}
An identical PPO agent that receives the full 16-dimensional state (Eq.~\ref{eq:state}), including the five disruption flags and five FTOPSIS closeness coefficients. By augmenting the observation with domain knowledge, M\textsubscript{3} can, in principle, exploit sustainability-resilience information from the first training step.

Both RL agents use identical hyperparameters to ensure a fair comparison: learning rate $3\!\times\!10^{-4}$, entropy coefficient $0.01$, clipping parameter $\epsilon=0.2$, mini-batch size~64, rollout length~2048, GAE~$\lambda=0.95$, 10~policy-update epochs per rollout, value-function coefficient~0.5, and maximum gradient norm~0.5.
Each agent is trained for $3\!\times\!10^6$~time steps across eight parallel environments (via \texttt{SubprocVecEnv}) with an evaluation callback every 160\,000~steps (5~evaluation episodes per checkpoint); the best checkpoint by mean evaluation reward is retained.
The extended training budget was chosen to ensure convergence of M\textsubscript{3}'s 16-dimensional observation space.

\subsubsection{Safety and Checkpoint Mechanisms}
\label{ssec:safety}
% Trimmed: The original two-layer description repeated the EvalCallback
% logic in full. Because Section 4.7 provides empirical evidence for
% the checkpoint mechanism, the methodology section now gives only a
% concise protocol summary, saving ~80 words.
Training for $3\!\times\!10^{6}$~timesteps risks late-stage policy collapse (catastrophic forgetting), a known phenomenon in on-policy algorithms~\citep{Schulman_etal_2017}.
We mitigate this with a two-layer safety protocol.
First, an \texttt{EvalCallback} evaluates the current policy every 160\,000~timesteps (five independent episodes); if the mean evaluation reward exceeds all previous checkpoints, the model is saved as \texttt{best\_model.zip}.
Second, periodic full-model snapshots are retained for post-hoc trajectory analysis and rollback.
Together, these mechanisms ensure that all reported metrics reflect the peak generalised policy, not the potentially degraded final state (see Section~\ref{sec:learning} for empirical validation).

\subsubsection{Experimental Design}
\label{ssec:exp_design}
To assess robustness across operating regimes, three demand-disruption scenarios of increasing severity are defined (Table~\ref{tab:scenarios}). Each of the nine model-scenario combinations is replicated with three independent random seeds (42, 123, 456), yielding 27~experimental runs.

\begin{table}[ht]
\centering
\caption{Scenario parameterisation. $\sigma_u$ denotes demand noise standard deviation (Product~1); $p_{\text{shock}}$ is daily demand-shock probability; $p_{\text{disrupt}}$ gives per-supplier daily disruption probabilities.}
\label{tab:scenarios}
\small
\begin{tabular}{@{}lccc@{}}
\toprule
Parameter & Stable Ops. & High Volatility & Systemic Shock \\
\midrule
$\sigma_1$ & 10 & 70 & 40 \\
$p_{\text{shock}}$ & 0.03 & 0.12 & 0.08 \\
Seasonality amplitude & 0.18 & 0.35 & 0.25 \\
$p_{\text{disrupt}}$ (S1--S5) & 0.05--0.03 & 0.18--0.12 & 0.25--0.05 \\
Systemic shock window & --- & --- & Days 50--250 \\
Failing suppliers & --- & --- & S1, S2, S4 \\
\bottomrule
\end{tabular}
\end{table}

The Systemic Shock scenario is deliberately \emph{truth-aligned}: the suppliers that fail (S1, S2, S4) have the lowest FTOPSIS resilience scores, while the survivors (S3, S5) rank highest. This design tests whether the MCDM-derived priors provide actionable value when the simulation physics mirror expert assessments.

% ==============================
% SECTION 4 --- RESULTS
% ==============================
%% ============================================================
%%  SECTION 4 — RESULTS AND DISCUSSION
%%  3 models × 3 scenarios × 3 seeds × 3 × 10^6 timesteps
%%  All statistics computed on 150 episodes per model–scenario
%% ============================================================

\section{Results and Discussion}
\label{sec:results}

This section reports the computational outcomes of the two-phase framework described in Section~\ref{sec:methodology}.
Phase~I (FBWM-FTOPSIS) produced a five-supplier sustainability-resilience ranking that remains fixed throughout the study (Table~\ref{tab:ftopsis_scores}).
Phase~II trained and evaluated three ordering policies, namely M\textsubscript{1}~(Base-Stock), M\textsubscript{2}~(Vanilla PPO), and M\textsubscript{3}~(PPO + FTOPSIS Priors), across three demand-disruption scenarios of increasing severity, each replicated with three independent random seeds ($s \in \{42, 123, 456\}$), for a total of 27~experimental runs.
Every reinforcement-learning run consumed $3\!\times\!10^{6}$~PPO timesteps distributed across eight parallel sub-environments, with 18~evaluation checkpoints (one every 160\,000~steps, each averaging five episodes); only the checkpoint with the highest mean evaluation reward was retained for post-training assessment (\texttt{EvalCallback} with \texttt{best\_model.zip}).
Post-training evaluation comprised 50~episodes per seed, yielding 150~episodes per model-scenario cell and 1\,350~episodes overall.
All experiments were executed on an Apple~M4 processor (10~CPU cores, 16\,GB unified memory); total wall-clock time was approximately 4.4~hours.

\begin{table}[htbp]
\centering
\caption{FTOPSIS closeness coefficients ($CC_i$) for the five-supplier pool from Phase~I.
Higher $CC_i$ indicates stronger sustainability-resilience performance.
These scores enter the M\textsubscript{3} observation vector (Eq.~\ref{eq:state}) and the sustainability reward component (Eq.~\ref{eq:reward}).}
\label{tab:ftopsis_scores}
\small
\begin{tabular}{@{}c l l c l@{}}
\toprule
Rank & Supplier & Archetype & $CC_i$ & Priority \\
\midrule
1 & $S_5$ & Premium Innovator     & 0.881 & High \\
2 & $S_3$ & Resilient Incumbent   & 0.796 & High \\
3 & $S_2$ & Green Specialist      & 0.691 & High \\
4 & $S_4$ & Balanced Generalist   & 0.494 & Medium \\
5 & $S_1$ & Low-Cost Leader       & 0.045 & Low \\
\bottomrule
\end{tabular}
\end{table}

The discussion proceeds through five complementary lenses: (i)~aggregate cost performance and statistical significance, (ii)~cost-sustainability Pareto trade-offs, (iii)~supplier allocation patterns, (iv)~bullwhip-effect dampening, and (v)~training convergence and the role of checkpoint selection.

%% ----------------------------------------------------------
\subsection{Aggregate Cost Performance}
\label{sec:cost_results}
%% ----------------------------------------------------------

Table~\ref{tab:main_results} reports the cross-seed mean and standard deviation of total cost, sustainability score, and bullwhip ratio for every model-scenario pair.
Figure~\ref{fig:performance} renders the same data as grouped bar charts with per-seed markers overlaid.

\begin{table}[htbp]
\centering
\caption{Cross-seed performance summary (mean $\pm$ SD, 150~episodes per cell).
Cost in \euro\,M; sustainability is the FTOPSIS-weighted allocation score; bullwhip is order-variance/demand-variance.
Bold marks the best model per scenario.}
\label{tab:main_results}
\small
\begin{tabular}{@{}l l r@{\,$\pm$\,}r r@{\,$\pm$\,}r r@{\,$\pm$\,}r@{}}
\toprule
Scenario & Model & \multicolumn{2}{c}{Cost (\euro\,M)} & \multicolumn{2}{c}{Sustainability} & \multicolumn{2}{c}{Bullwhip} \\
\midrule
\multirow{3}{*}{Stable Ops.}
  & M\textsubscript{1} (Base-Stock)  & 45.27 & 0.55 & 0.573 & 0.000 & \textbf{0.00} & 0.00 \\
  & M\textsubscript{2} (Vanilla PPO) & 20.79 & 3.62 & \textbf{0.607} & 0.040 & 15.16 & 4.72 \\
  & M\textsubscript{3} (PPO + Priors) & \textbf{16.38} & 1.52 & 0.553 & 0.119 & 18.86 & 6.86 \\
\midrule
\multirow{3}{*}{High Vol.}
  & M\textsubscript{1} (Base-Stock)  & 34.43 & 1.01 & 0.573 & 0.000 & \textbf{0.00} & 0.00 \\
  & M\textsubscript{2} (Vanilla PPO) & \textbf{10.94} & 1.21 & \textbf{0.635} & 0.045 & 0.39 & 0.15 \\
  & M\textsubscript{3} (PPO + Priors) & 13.00 & 4.14 & 0.578 & 0.045 & 0.41 & 0.14 \\
\midrule
\multirow{3}{*}{Syst.\ Shock}
  & M\textsubscript{1} (Base-Stock)  & 19.05 & 0.62 & 0.573 & 0.000 & \textbf{0.00} & 0.00 \\
  & M\textsubscript{2} (Vanilla PPO) & \textbf{3.80} & 0.55 & 0.767 & 0.121 & 7.71 & 2.20 \\
  & M\textsubscript{3} (PPO + Priors) & 5.15 & 2.35 & \textbf{0.842} & 0.017 & 9.70 & 2.94 \\
\bottomrule
\end{tabular}
\end{table}

\begin{figure}[htbp]
    \centering
    \includegraphics[width=\textwidth]{figures/figure1_performance_comparison.png}
    \caption{Grouped bar chart of total cost, sustainability score, and bullwhip ratio across three scenarios.
    Bars show cross-seed means; scatter markers expose per-seed dispersion.
    Both RL agents compress costs by 54--80\,\% relative to M\textsubscript{1}.}
    \label{fig:performance}
\end{figure}

\subsubsection{RL versus the Deterministic Baseline}

The most prominent finding in Table~\ref{tab:main_results} is the magnitude of cost reduction that both PPO agents achieve relative to the Base-Stock heuristic.
Under Stable Operations, M\textsubscript{3} reduces total cost by 63.8\,\% (\euro\,16.38M versus \euro\,45.27M) and M\textsubscript{2} by 54.1\,\% (\euro\,20.79M versus \euro\,45.27M).
Under High Volatility, M\textsubscript{2} achieves a 68.2\,\% reduction (\euro\,10.94M versus \euro\,34.43M).
The largest single-scenario gap appears under Systemic Shock, where M\textsubscript{2} reaches \euro\,3.80M, an 80.1\,\% reduction from the \euro\,19.05M Base-Stock benchmark.
One-way ANOVA across all three models confirms that these differences are not artefacts of sample variability: $F = 6{,}924$ (Stable), $F = 3{,}881$ (High Volatility), and $F = 5{,}151$ (Systemic Shock), all with $p \approx 0$.
These results corroborate the emerging consensus in the RL-for-supply-chain literature that policy-gradient methods can discover dynamic ordering strategies fundamentally inaccessible to static heuristics~\citep{Wang_etal_2022,Li_etal_2024}.

\subsubsection{\texorpdfstring{M\textsubscript{3} versus M\textsubscript{2}}{M3 versus M2}: A Scenario-Dependent Trade-off}

The comparison between the two PPO agents, the central empirical question of this study, reveals a nuanced, scenario-dependent pattern rather than a uniform winner.

\paragraph{Stable Operations: peacetime efficiency.}
Under low noise and no disruptions, M\textsubscript{3} achieves a mean cost of \euro\,16.38M $\pm$ \euro\,1.52M, 21.2\,\% lower than M\textsubscript{2}'s \euro\,20.79M $\pm$ \euro\,3.62M.
A Welch $t$-test confirms this gap at $t = 13.75$, $p = 1.0\!\times\!10^{-30}$, with a large effect size of $d = 1.59$ (Table~\ref{tab:significance}).
The result is noteworthy because the FTOPSIS priors carry little marginal information in a benign environment: there are no disruptions to predict and no unreliable suppliers to avoid.
That M\textsubscript{3} still outperforms suggests the priors function as an implicit \emph{regulariser}: by structuring the 16-dimensional observation around supplier quality scores, they constrain the policy search to a lower-dimensional manifold that PPO traverses more efficiently than the unconstrained 6-dimensional space of M\textsubscript{2}.
However, M\textsubscript{2} leads on sustainability (0.607 vs.\ 0.553, $d = 0.62$, medium effect) and bullwhip ratio (15.16 vs.\ 18.86, $d = 0.58$, medium effect), indicating that M\textsubscript{3}'s cost advantage comes at the expense of order smoothness.

\paragraph{High Volatility: the complexity tax.}
When demand variance increases but supplier capacity remains intact, M\textsubscript{2} reclaims the cost advantage: \euro\,10.94M $\pm$ \euro\,1.21M versus M\textsubscript{3}'s \euro\,13.00M $\pm$ \euro\,4.14M ($t = -5.84$, $p = 2.6\!\times\!10^{-8}$, $d = 0.67$, medium).
M\textsubscript{2} also leads on sustainability (0.635 vs.\ 0.578, $d = 1.29$, large).
The bullwhip ratios, however, are statistically indistinguishable: 0.39 versus 0.41 ($p = 0.15$, $d = 0.17$, negligible).
The explanation lies in the nature of High Volatility: demand shocks affect all suppliers equally, so the disruption flags in M\textsubscript{3}'s observation carry no discriminative signal, while the 10~additional input dimensions impose a learning burden.
The standard deviation of M\textsubscript{3}'s cost ($\pm$\,\euro\,4.14M, three times that of M\textsubscript{2}) reveals that one of the three seeds (seed~456, \euro\,17.72M) converges to a substantially weaker policy than the other two (seed~42, \euro\,9.50M; seed~123, \euro\,11.77M), exposing the sensitivity of the 16-dimensional agent to random initialisation when the environment does not fully activate its extra channels.

\paragraph{Systemic Shock: wartime resilience.}
The most demanding scenario reverses the narrative on sustainability.
M\textsubscript{3} achieves a sustainability score of 0.842 $\pm$ 0.017, a 9.8\,\% improvement over M\textsubscript{2}'s 0.767 $\pm$ 0.121 ($t = -7.49$, $p = 4.9\!\times\!10^{-12}$, $d = 0.86$, large).
The standard deviation tells an equally important story: M\textsubscript{3}'s sustainability varies by only $\pm$\,0.017 across seeds, whereas M\textsubscript{2}'s $\pm$\,0.121 reflects a seven-fold wider spread.
M\textsubscript{2} leads on cost (\euro\,3.80M vs.\ \euro\,5.15M, $d = 0.79$, medium), but this advantage is driven primarily by seed~456, in which M\textsubscript{2} achieves a low sustainability of 0.596 (far below the other two seeds' $>$\,0.85), suggesting the agent discovered a cost-minimising but sustainability-poor policy.
M\textsubscript{3} thus provides a more \emph{predictable} and \emph{sustainability-aligned} response to systemic disruptions, a property of direct value to risk-averse supply chain managers.

%% ----------------------------------------------------------
\subsection{Statistical Significance}
\label{sec:statistics}
%% ----------------------------------------------------------

Table~\ref{tab:significance} formalises the pairwise comparisons between M\textsubscript{2} and M\textsubscript{3} using Welch's unequal-variance $t$-tests (150~episodes per group) and Cohen's $d$ effect sizes.
Of the nine metric-scenario cells, eight reach conventional significance ($p < 0.05$); the sole exception is the bullwhip ratio under High Volatility ($p = 0.15$, $d = 0.17$), where both agents achieve near-identical demand-smoothing behaviour with ratios below unity.

\begin{table}[htbp]
\centering
\caption{Pairwise tests: M\textsubscript{2} vs.\ M\textsubscript{3} (Welch's $t$, two-tailed, $\alpha = 0.05$).
``Favoured'' indicates the superior model.
Effect sizes: $|d| < 0.2$ negligible, $0.2$--$0.5$ small, $0.5$--$0.8$ medium, $> 0.8$ large; bold $d$ denotes large effects.}
\label{tab:significance}
\small
\begin{tabular}{@{}l l r r r l@{}}
\toprule
Scenario & Metric & $t$ & $p$ & $|d|$ & Favoured \\
\midrule
\multirow{3}{*}{Stable Ops.}
  & Cost        & $+13.75$ & $1.0\!\times\!10^{-30}$ & \textbf{1.59} & M\textsubscript{3} \\
  & Sustainability & $+5.33$ & $3.0\!\times\!10^{-7}$ & 0.62 & M\textsubscript{2} \\
  & Bullwhip    & $-5.06$  & $8.0\!\times\!10^{-7}$  & 0.58 & M\textsubscript{2} \\
\midrule
\multirow{3}{*}{High Vol.}
  & Cost        & $-5.84$  & $2.6\!\times\!10^{-8}$  & 0.67 & M\textsubscript{2} \\
  & Sustainability & $+11.16$ & $2.3\!\times\!10^{-24}$ & \textbf{1.29} & M\textsubscript{2} \\
  & Bullwhip    & $-1.44$  & $0.151$                 & 0.17 & --- \\
\midrule
\multirow{3}{*}{Syst.\ Shock}
  & Cost        & $-6.87$  & $1.3\!\times\!10^{-10}$ & 0.79 & M\textsubscript{2} \\
  & Sustainability & $-7.49$ & $4.9\!\times\!10^{-12}$ & \textbf{0.86} & M\textsubscript{3} \\
  & Bullwhip    & $-6.34$  & $9.5\!\times\!10^{-10}$ & 0.73 & M\textsubscript{2} \\
\bottomrule
\end{tabular}
\end{table}

Two large-effect findings anchor the remainder of the discussion:
(i)~M\textsubscript{3} reduces cost in Stable Operations with $d = 1.59$, the single largest effect in the study, and
(ii)~M\textsubscript{3} raises sustainability under Systemic Shock with $d = 0.86$, the only scenario where M\textsubscript{3} achieves a statistically significant advantage on a metric other than cost.
Together, these two results delineate the operational envelope of the MCDM-augmented agent: it excels at \emph{cost compression} when the environment is calm and at \emph{sustainability assurance} when the environment is hostile.

%% ----------------------------------------------------------
\subsection{Seed Variability and Robustness}
\label{sec:variability}
%% ----------------------------------------------------------

Multi-seed replication exposes important differences in learning robustness that single-seed experiments would mask.
Figure~\ref{fig:variability} displays the per-seed cost distributions for M\textsubscript{2} and M\textsubscript{3}.

\begin{figure}[htbp]
    \centering
    \includegraphics[width=\textwidth]{figures/figure2_seed_variability.png}
    \caption{Seed-level cost variability for M\textsubscript{2} and M\textsubscript{3}.
    Each strip shows the 50-episode distribution for one seed; box plots overlay the interquartile range.
    M\textsubscript{3} clusters tightly under Stable Operations (spread ratio~1.2$\times$) but diverges under High Volatility (seed~456 at \euro\,17.7M vs.\ ${\sim}$\euro\,10M for other seeds).}
    \label{fig:variability}
\end{figure}

Under Stable Operations, the three M\textsubscript{3} seeds converge to a narrow cost band (\euro\,15.31M, \euro\,15.34M, \euro\,18.48M; inter-seed spread ratio~1.2$\times$), tighter than M\textsubscript{2}'s range (\euro\,16.70M, \euro\,20.20M, \euro\,25.46M; spread ratio~1.5$\times$).
This is a marked improvement over the $1\!\times\!10^{6}$-step training regime, where M\textsubscript{3} under Stable Operations exhibited an 18$\times$ spread ratio with one seed trapped in a poor local optimum.
The extended $3\!\times\!10^{6}$-step budget allowed all three M\textsubscript{3} seeds to escape early attractors and converge to similarly cost-effective policies, confirming that the 16-dimensional observation space requires a larger training investment but rewards it with stable convergence.

Under High Volatility, M\textsubscript{3} shows the reverse pattern: seed~456 converges to a substantially weaker policy (\euro\,17.72M versus \euro\,9.50M and \euro\,11.77M for the other seeds).
Because demand shocks affect all suppliers symmetrically in this scenario, the disruption flags and FTOPSIS scores carry limited discriminative value, and the expanded observation space occasionally guides a seed into a suboptimal basin.
Under Systemic Shock, M\textsubscript{3}'s seed~123 (\euro\,8.33M) underperforms the other two seeds (\euro\,3.78M, \euro\,3.34M), though even the worst-case M\textsubscript{3} seed still reduces cost by 56\,\% relative to M\textsubscript{1}.

These patterns suggest that the MCDM priors are most stabilising when the environment dynamics align with the information they encode: the priors compress the inter-seed variance when disruptions discriminate among suppliers (Stable Operations), but inflate it when the disruption structure is diffuse (High Volatility).

%% ----------------------------------------------------------
\subsection{Cost-Sustainability Pareto Analysis}
\label{sec:pareto}
%% ----------------------------------------------------------

\begin{figure}[htbp]
    \centering
    \includegraphics[width=0.85\textwidth]{figures/figure3_pareto_cost_sustainability.png}
    \caption{Cost-sustainability Pareto frontier.
    Each marker is a cross-seed mean; error bars show $\pm$\,1~SD.
    M\textsubscript{1} is Pareto-dominated in every scenario.
    Under Systemic Shock, M\textsubscript{3} achieves the study's highest sustainability (0.842) at a modest cost premium over M\textsubscript{2}.}
    \label{fig:pareto}
\end{figure}

Figure~\ref{fig:pareto} projects every model-scenario combination into the two-dimensional cost-sustainability plane.
M\textsubscript{1} is Pareto-dominated in all three scenarios, reinforcing the case for RL-based policies.
Among the RL agents, the Pareto-efficient set reveals a scenario-dependent complementarity.
M\textsubscript{2} under Systemic Shock occupies the lowest-cost position on the frontier (\euro\,3.80M, sustainability~0.767), while M\textsubscript{3} under the same scenario achieves the highest sustainability of any cell in the study (0.842) at a cost of \euro\,5.15M.
A supply chain manager who values sustainability assurance under disruption, as many firms increasingly do under ESG mandates~\citep{Ivanov_Dolgui_2021}, would rationally select M\textsubscript{3}'s Systemic Shock policy as the preferred operating point, accepting a \euro\,1.35M cost premium for a 9.8\,\% sustainability gain and a seven-fold reduction in cross-seed sustainability variance.

Under Stable Operations, M\textsubscript{3} dominates M\textsubscript{2} on cost (\euro\,16.38M vs.\ \euro\,20.79M) but cedes ground on sustainability (0.553 vs.\ 0.607).
Neither agent dominates the other in the Pareto sense, indicating that the choice between them depends on the manager's relative weighting of cost versus sustainability.
Under High Volatility, M\textsubscript{2} dominates M\textsubscript{3} on both axes, making it the clear selection for that regime.

%% ----------------------------------------------------------
\subsection{Supplier Allocation Patterns}
\label{sec:allocation}
%% ----------------------------------------------------------

\begin{figure}[htbp]
    \centering
    \includegraphics[width=\textwidth]{figures/figure4_allocation_heatmap.png}
    \caption{Mean daily order quantities allocated to each supplier by model and scenario.
    Under Systemic Shock, M\textsubscript{3} routes 77.8\,\% of volume to $S_5$ versus 59.1\,\% for M\textsubscript{2}; both agents virtually eliminate $S_1$.}
    \label{fig:allocation}
\end{figure}

The supplier allocation heatmap (Figure~\ref{fig:allocation}) reveals how each agent distributes orders across the five-supplier pool and whether those patterns align with the Phase~I FTOPSIS ranking.

\paragraph{Systemic Shock: FTOPSIS-aligned concentration.}
The Systemic Shock scenario was designed with a truth-aligned disruption structure: the suppliers most likely to fail (S\textsubscript{1}, S\textsubscript{2}, S\textsubscript{4}) are precisely those with lower FTOPSIS scores, while the survivors (S\textsubscript{3}, S\textsubscript{5}) rank highest.
Both RL agents learn to exploit this structure, but to strikingly different degrees.
M\textsubscript{3} directs 77.8\,\% of its 136.3~daily units to $S_5$ (the top-FTOPSIS supplier) and 17.4\,\% to $S_3$ (the second-ranked), channelling 95.2\,\% of total volume to the two most resilient suppliers.
M\textsubscript{2}, without explicit access to the FTOPSIS scores, nonetheless routes 59.1\,\% of its 155.8~daily units to $S_5$ and 19.9\,\% to $S_3$, reaching a combined 79.0\,\%.
This convergence of both agents toward the FTOPSIS-recommended ranking, despite M\textsubscript{2} having no direct sustainability signal, provides strong indirect evidence that the simulation's disruption dynamics encode the same quality that the MCDM evaluation captures: suppliers with higher sustainability-resilience scores are also more likely to remain operational under stress.

The 18.7-percentage-point gap in $S_5$ concentration (77.8\,\% vs.\ 59.1\,\%) is the mechanism behind M\textsubscript{3}'s sustainability advantage ($d = 0.86$).
By frontloading the FTOPSIS scores into the observation, M\textsubscript{3} does not need to discover the supplier ranking through trial and error; it can allocate to $S_5$ from the first training step and refine only the quantity.

\paragraph{Stable Operations: comparable volume, different supplier mix.}
Under benign conditions, both RL agents place nearly identical total daily orders (M\textsubscript{2}: 161.0~units; M\textsubscript{3}: 158.7~units), far below M\textsubscript{1}'s 500.8~units.
The cost difference between M\textsubscript{2} and M\textsubscript{3} therefore derives not from volume but from unit-cost composition: M\textsubscript{3} allocates more to $S_1$ (33.9~units vs.\ 18.2~units), the cheapest but lowest-FTOPSIS supplier, which reduces total expenditure at the expense of the sustainability score.

\paragraph{High Volatility: diversification without discrimination.}
Under High Volatility, M\textsubscript{2} spreads orders across $S_2$--$S_3$ (49.1 and 58.0~units, together 58.7\,\% of 182.3~total), while M\textsubscript{3} distributes more evenly, with no single supplier exceeding 27.1\,\% of its 194.7~daily total.
Neither agent achieves a clear FTOPSIS-aligned ranking under this scenario, consistent with the observation that symmetric demand shocks provide no supplier-discriminative learning signal.

%% ----------------------------------------------------------
\subsection{Bullwhip Effect}
\label{sec:bullwhip}
%% ----------------------------------------------------------

\begin{figure}[htbp]
    \centering
    \includegraphics[width=0.75\textwidth]{figures/figure5_bullwhip_comparison.png}
    \caption{Bullwhip ratio for M\textsubscript{2} and M\textsubscript{3} across three scenarios.
    M\textsubscript{1} is omitted (bullwhip is zero by construction).
    Under High Volatility, both agents achieve ratios below unity, actively dampening demand variability.}
    \label{fig:bullwhip}
\end{figure}

The bullwhip effect, the amplification of demand variability as orders propagate upstream, is a classical concern in supply chain management~\citep{Ivanov_Dolgui_2021}.
Figure~\ref{fig:bullwhip} compares the bullwhip ratios of M\textsubscript{2} and M\textsubscript{3} across all three scenarios.

\paragraph{Demand dampening under High Volatility.}
Perhaps the most striking finding in the bullwhip analysis is that both RL agents achieve ratios \emph{below unity} under High Volatility: M\textsubscript{2} at 0.39 and M\textsubscript{3} at 0.41.
A bullwhip ratio below 1.0 means the agent's order stream is \emph{smoother} than the incoming demand signal; the agent absorbs demand shocks through inventory buffers rather than passing them upstream.
This is the opposite of the demand-amplification pathology that plagues most inventory systems and provides direct evidence that PPO discovers order-smoothing strategies without explicit instruction.
The difference between the two agents is statistically negligible ($p = 0.15$, $d = 0.17$), indicating that the MCDM priors neither help nor hinder the dampening mechanism in this regime.

\paragraph{Amplification under Stable Operations and Systemic Shock.}
Under Stable Operations, both agents produce bullwhip ratios substantially above unity (M\textsubscript{2}: 15.16; M\textsubscript{3}: 18.86), indicating order amplification.
This is a known artefact of low-variance demand: when demand is nearly constant, even modest order fluctuations produce a high variance \emph{ratio}, though the absolute order variance may be small.
Under Systemic Shock, M\textsubscript{2} amplifies less than M\textsubscript{3} (7.71 vs.\ 9.70, $d = 0.73$, medium), reflecting M\textsubscript{2}'s more diversified allocation across suppliers (total 155.8~units vs.\ M\textsubscript{3}'s 136.3~units), which mechanically distributes variance across more channels.

%% ----------------------------------------------------------
\subsection{Training Convergence and Checkpoint Selection}
\label{sec:learning}
%% ----------------------------------------------------------

\begin{figure}[htbp]
    \centering
    \includegraphics[width=\textwidth]{figures/figure7_learning_curves.png}
    \caption{Evaluation reward trajectories over the $3\!\times\!10^{6}$-step training horizon (18 checkpoints).
    Solid lines: cross-seed mean; shaded bands: $\pm$\,1~SD.
    Panel~(a) shows a post-peak collapse under Stable Operations, validating the \texttt{EvalCallback} safety mechanism.
    Panel~(c) shows M\textsubscript{3} still improving at the final checkpoint under Systemic Shock.}
    \label{fig:learning}
\end{figure}

Figure~\ref{fig:learning} plots the evaluation reward trajectories for all 18~RL training runs across the full $3\!\times\!10^{6}$-step horizon.
We deliberately present the complete, untruncated curves so that the late-stage behaviour is fully transparent.
Three patterns merit discussion.

\paragraph{Late-stage policy collapse and the role of \texttt{EvalCallback}.}
Under Stable Operations, M\textsubscript{3}'s cross-seed mean reward peaks near step~2.4M and then collapses; in the most extreme run (seed~42), the peak reward of $-1{,}512$ degrades to $-4{,}813$ at step~2.88M (218\,\% decline).
This is a textbook manifestation of catastrophic forgetting in on-policy algorithms~\citep{Schulman_etal_2017}.
The \texttt{EvalCallback} safety mechanism (Section~\ref{ssec:safety}) ensures that reported metrics reflect the peak checkpoint: had we used the final checkpoint, M\textsubscript{3}'s Stable Operations cost would have been ${\sim}$\euro\,51M instead of \euro\,16.38M.

\paragraph{Systemic Shock: late convergence of M\textsubscript{3}.}
Under Systemic Shock, all three M\textsubscript{3} seeds achieve their peak evaluation reward at the final checkpoint (step~2.88M), with zero reward divergence between peak and terminal evaluations.
This late convergence stands in sharp contrast to M\textsubscript{2}, which peaks between steps~1.4M and~2.6M and then drifts by $-29$ to $-38$ reward units.
The implication is that M\textsubscript{3}'s 16-dimensional observation space requires more gradient steps to fully exploit the FTOPSIS priors under disruption conditions, and may yet improve with additional training beyond $3\!\times\!10^{6}$~steps.
This pattern reinforces the interpretation of MCDM priors as a deferred asset: they impose a higher initial learning cost but continue to yield returns at late training stages where simpler observation spaces have already plateaued.

\paragraph{Convergence rate in the first million steps.}
Under all three scenarios, both agents achieve the majority of their improvement within the first $1.0\!\times\!10^{6}$~steps, with diminishing returns thereafter.
The extended $3\!\times\!10^{6}$-step budget was most valuable for M\textsubscript{3} under Stable Operations, where all three seeds converged to a tight cost band (\euro\,15.3M--\euro\,18.5M) after requiring $>$\,1.7M steps to escape early local optima, a substantial improvement over the 18$\times$ inter-seed spread ratio observed in preliminary $1\!\times\!10^{6}$-step experiments.

%% ----------------------------------------------------------
\subsection{Discussion and Synthesis}
\label{sec:synthesis}
%% ----------------------------------------------------------

The results cohere into three empirically grounded propositions.

\paragraph{Proposition 1: RL dominates static heuristics across all operating regimes.}
Both PPO agents reduce total cost by 54--80\,\% relative to the Base-Stock heuristic in every scenario, with ANOVA $F$-statistics ranging from 3{,}881 to 6{,}924 ($p \approx 0$).
This advantage is not merely statistical: the smallest absolute gap (\euro\,13.90M under Systemic Shock, M\textsubscript{2} vs.\ M\textsubscript{1}) represents a 73\,\% reduction.
The result is consistent with prior findings in multi-echelon inventory RL~\citep{Wang_etal_2022,Li_etal_2024} and extends them to a multi-supplier, multi-objective setting with supplier disruptions.

\paragraph{Proposition 2: MCDM priors are a regime-activated asset, not a universal upgrade.}
M\textsubscript{3}'s value depends on whether the environment activates the information embedded in its observation space.
Under Stable Operations, where no disruptions occur, the priors still confer a cost advantage ($d = 1.59$) by regularising the policy search.
Under Systemic Shock, where the disruption structure mirrors the FTOPSIS ranking, the priors yield a large sustainability advantage ($d = 0.86$) and near-perfect alignment of orders with the top-FTOPSIS supplier (77.8\,\% of volume to $S_5$).
Under High Volatility, where demand shocks are supplier-agnostic, the priors add dimensionality without discriminative value, and M\textsubscript{2} outperforms.
This three-regime pattern is consistent with Bayesian model-selection theory: informative priors improve estimation under structured uncertainty but can introduce bias under diffuse noise~\citep{Chen_Kumar_2024}.
The practical implication is that MCDM-prior augmentation should be deployed selectively, targeted at environments where expert domain knowledge is expected to align with the true disruption structure.

\paragraph{Proposition 3: checkpoint selection is a safety-critical mechanism.}
The divergence between peak and terminal evaluation rewards (up to 218\,\% degradation in the worst case) demonstrates that on-policy RL agents cannot be deployed from their final training state.
The \texttt{EvalCallback} with \texttt{best\_model.zip} retention functions as a safeguard against late-stage policy collapse, analogous to early stopping in supervised learning.
Future deployments of RL-based ordering systems should treat checkpoint selection as a first-class safety mechanism, with evaluation protocols calibrated to the decision-relevant metrics (cost, sustainability, bullwhip) rather than raw reward alone.

\paragraph{Managerial implications.}
The results suggest a conditional deployment strategy.
In stable, predictable environments, M\textsubscript{3} achieves the lowest cost (\euro\,16.38M, a 63.8\,\% saving over the heuristic baseline) and should be preferred when cost minimisation is the primary objective.
When demand volatility is high but supplier networks remain intact, M\textsubscript{2}'s simpler architecture generalises more reliably and should be selected.
When systemic disruptions are anticipated, the scenario of greatest managerial concern, M\textsubscript{3} delivers the highest sustainability (0.842), the tightest sustainability variance ($\pm$\,0.017), and routes 95\,\% of orders to the two most resilient suppliers.
A risk-averse manager facing ESG mandates or regulatory sustainability requirements would rationally accept M\textsubscript{3}'s \euro\,1.35M cost premium (a 35.5\,\% increase over M\textsubscript{2}, but still 73.0\,\% below the heuristic) for this level of sustainability assurance.

\paragraph{Limitations.}
Four limitations merit transparent acknowledgment.
First, the study uses three random seeds per configuration, which is standard in deep RL benchmarking~\citep{Schulman_etal_2017} but limits the precision of variability estimates; five or more seeds would tighten the confidence intervals on inter-seed spread.
Second, the late convergence of M\textsubscript{3} under Systemic Shock (peak at the final checkpoint) suggests that additional training beyond $3\!\times\!10^{6}$~steps could further improve M\textsubscript{3}'s cost performance in that regime; the current results may underestimate M\textsubscript{3}'s ceiling.
Third, the truth-aligned disruption structure of the Systemic Shock scenario (failing suppliers have low FTOPSIS scores) represents a best-case test for the MCDM priors; a misaligned scenario, where the highest-rated supplier fails, would provide a sterner test of robustness.
Fourth, sensitivity analysis over the sustainability weight $\lambda_{\text{sust}}$ has not been conducted for the $3\!\times\!10^{6}$-step regime; the interaction between training budget and reward weighting is an open question for future work.

% ==============================
% SECTION 5 --- CONCLUSION
% ==============================
\section{Conclusion}
\label{sec:conclusion}

This study introduced a two-phase hybrid framework that integrates Fuzzy Best-Worst Method and Fuzzy TOPSIS (FBWM-FTOPSIS) with Proximal Policy Optimization (PPO) for dynamic supplier order allocation under uncertainty.
Phase~I translated expert sustainability and resilience judgments across 16~criteria into quantitative closeness coefficients for five heterogeneous suppliers.
Phase~II embedded these coefficients directly into the PPO agent's observation vector, creating a prior-augmented agent (M\textsubscript{3}) that was benchmarked against a vanilla PPO baseline (M\textsubscript{2}) and a deterministic base-stock heuristic (M\textsubscript{1}).
All 27~experimental configurations (three models, three demand-disruption scenarios, three random seeds) were trained for $3\!\times\!10^{6}$~timesteps and evaluated over 1\,350~episodes.
The findings cohere into three core propositions.

% Trimmed: The original three proposition paragraphs repeated the detailed
% statistics already presented in Section 4.8 (Discussion and Synthesis).
% The revision distils each proposition to its core claim with minimal
% duplication, saving approximately 200 words.

\paragraph{Proposition 1: RL dominates static heuristics unconditionally.}
Both PPO agents reduced total supply-chain cost by 54--80\,\% relative to the Base-Stock heuristic across every scenario (ANOVA $F = 3{,}881$--$6{,}924$, $p \approx 0$), confirming that policy-gradient methods discover dynamic ordering strategies inaccessible to static rules and extending prior single-objective findings~\citep{Wang_etal_2022,Li_etal_2024} to a multi-supplier, multi-objective setting with correlated disruptions.

\paragraph{Proposition 2: MCDM priors are a regime-activated asset.}
The value of domain-expert priors depends on whether the operating environment activates the information they encode.
M\textsubscript{3} delivered a large cost advantage under Stable Operations ($d = 1.59$) and the highest sustainability under Systemic Shock ($d = 0.86$, with seven-fold tighter cross-seed variance), but incurred a complexity tax under High Volatility, where demand shocks are supplier-agnostic.
This pattern is consistent with Bayesian model-selection theory: informative priors accelerate learning under structured uncertainty but can introduce bias under diffuse noise~\citep{Chen_Kumar_2024}.

\paragraph{Proposition 3: RL agents can absorb, not merely transmit, demand shocks.}
Both agents achieved bullwhip ratios below unity under High Volatility (0.39 and 0.41), actively dampening demand variability rather than amplifying it. Because the dampening was statistically indistinguishable between agents ($p = 0.15$), it appears to arise from the RL mechanism itself rather than from the MCDM priors.

\paragraph{Practical implications.}
For supply chain managers, the results suggest a conditional deployment strategy.
In stable, predictable environments, M\textsubscript{3} should be preferred for its cost advantage (63.8\,\% saving over the heuristic).
When systemic disruptions are anticipated, M\textsubscript{3} provides the highest sustainability assurance (0.842) and the most predictable behaviour across seeds, properties valued by risk-averse firms operating under ESG mandates.
When demand volatility is high but supplier networks remain intact, M\textsubscript{2}'s simpler architecture generalises more reliably and should be selected.
Collectively, these capabilities position the framework as a tool for cleaner and more sustainable procurement in industries where supplier environmental and social performance directly affects the ecological footprint of the buying organisation.

\paragraph{Methodological contribution.}
Beyond the empirical findings, this study makes a methodological contribution by demonstrating that \texttt{EvalCallback}-based checkpoint selection is a safety-critical component of RL-based supply chain systems.
The divergence between peak and terminal evaluation rewards reached 218\,\% in the worst case; had we deployed the final checkpoint rather than the best, M\textsubscript{3}'s cost under Stable Operations would have been approximately \euro\,51M instead of \euro\,16.38M.
This finding underscores that checkpoint selection must be treated as a first-class safety mechanism in any operational deployment of RL-based ordering.

\paragraph{Limitations and future work.}
The study's main limitations (three seeds per configuration, truth-aligned disruption structure, late convergence under Systemic Shock, and unexplored $\lambda_{\text{sust}}$ sensitivity) are detailed in Section~\ref{sec:synthesis}.
The most promising direction is an \emph{adaptive observation-masking} mechanism that dynamically activates or deactivates the MCDM channels based on a real-time volatility estimate, potentially yielding a single agent that dominates across all regimes.

% ==============================
% BACK-MATTER SECTIONS (Phase 2)
% ==============================
% JCP requires these sections in this exact order, placed after
% the main text and before the reference list.

\section*{CRediT authorship contribution statement}

\textbf{Ali Vaezi:} Conceptualization, Methodology, Software, Formal analysis,
Investigation, Data curation, Visualization, Writing -- original draft,
Writing -- review \& editing, Project administration.
\textbf{Erfan Rabbani:} Conceptualization, Methodology, Supervision,
Writing -- review \& editing.

\section*{Declaration of competing interest}

The authors declare that they have no known competing financial interests
or personal relationships that could have appeared to influence the work
reported in this paper.

\section*{Funding}

This research did not receive any specific grant from funding agencies
in the public, commercial, or not-for-profit sectors.

\section*{Data availability}

% OPTION A (preferred): Uncomment the following two lines after depositing
% code and data on Zenodo, replacing XXXXXXX with the actual DOI:
% The source code, trained model checkpoints, and experimental result files
% supporting the findings of this study are publicly available at
% \url{https://doi.org/10.5281/zenodo.XXXXXXX}.
%
% OPTION B (current fallback):
Data will be made available on request.

\section*{Declaration of generative AI and AI-assisted technologies in the manuscript preparation process}

% NOTE: Remove this entire section if the authors decide that no
% generative AI was used at any stage of the research or writing.
During the preparation of this work the authors used GitHub Copilot
in order to assist with Python code development for the reinforcement
learning experiments and with \LaTeX{} formatting of the manuscript.
After using this tool, the authors reviewed and edited all content
as needed and take full responsibility for the content of the
published article.

% Bibliography style changed from unsrtnat (numbered) to elsarticle-harv
% (author-year Harvard). This is required because JCP uses an author-year
% reference format. The elsarticle-harv style produces citations such as
% "(Ivanov and Dolgui, 2021)" and alphabetically sorted reference lists,
% matching the journal's published papers.
\bibliographystyle{elsarticle-harv}
\begingroup
  \raggedright            % Allow ragged-right setting inside the reference list
  \emergencystretch=2em   % Extra flexibility to suppress underfull \hbox from long DOIs
  \bibliography{bibliography}
\endgroup

\end{document}
